\section{Telaah Beberapa Cipher Blok Simetris Lainnya}

Selain DES dan AES, komunitas kriptografi internasional telah mengembangkan berbagai algoritma cipher blok yang memiliki keunikan arsitektur, parameter fleksibel, dan filosofi perancangan tersendiri. Bagian ini mengkaji algoritma standar nasional Rusia (GOST), cipher inovatif berbasis rotasi dan data-dependent (RC5 dan RC6), serta deretan algoritma finalis kompetisi AES NIST (Blowfish, Twofish, Serpent, dan MARS), sebelum ditutup dengan telaah penerapannya dalam infrastruktur kehidupan modern.

\subsection{Algoritma GOST 28147-89 (Magma)}

\textbf{GOST 28147-89} (kini diperbarui dalam standar GOST R 34.12-2015 dengan nama sandi \textbf{Magma}) adalah standar algoritma enkripsi simetris nasional Uni Soviet yang pertama kali dirilis secara rahasia pada tahun 1970-an dan dideklasifikasi untuk umum pada tahun 1989. GOST berfungsi sebagai padanan resmi Soviet terhadap standar DES Amerika Serikat.

\begin{figure}[htbp]
  \centering
  \includegraphics[width=0.75\textwidth]{bab-05/gost-magma-putaran}
  \caption{Struktur satu putaran jaringan Feistel pada algoritma GOST 28147-89 (Magma).}
  \label{fig:gost-putaran}
\end{figure}

Karakteristik arsitektur GOST:
\begin{itemize}
  \item \textbf{Ukuran Blok}: 64 bit.
  \item \textbf{Panjang Kunci}: 256 bit—panjang kunci yang sangat masif dan visioner untuk era 1970-an (jauh melampaui kunci 56-bit DES).
  \item \textbf{Struktur}: Jaringan Feistel dengan \textbf{32 putaran} (dua kali lipat jumlah putaran DES).
\end{itemize}

\subsubsection{Fungsi Putaran dan S-box GOST}

Pada setiap putaran ke-$i$, paruh kanan 32-bit $R_{i-1}$ diproses bersama subkunci 32-bit $K_i$ melalui langkah-langkah:
\begin{enumerate}
  \item \textbf{Penjumlahan Modular}: $A = (R_{i-1} + K_i) \bmod 2^{32}$. Berbeda dengan DES yang menggunakan XOR, GOST memanfaatkan penjumlahan aritmatika modulo $2^{32}$.
  \item \textbf{Substitusi Kotak-S}: Vektor 32 bit $A$ dibagi menjadi delapan kelompok 4-bit ($A_0, \dots, A_7$). Masing-masing dimasukkan ke dalam 8 buah kotak-S $4 \times 4$ ($S_0, \dots, S_7$) yang menghasilkan keluaran 4 bit.
  \item \textbf{Rotasi Siklis}: Rangkaian 32 bit keluaran kotak-S diputar secara siklis ke kiri sejauh \textbf{11 bit} ($X \lll 11$) untuk menghasilkan difusi yang cepat antar-kotak.
\end{enumerate}

\begin{figure}[htbp]
  \centering
  \includegraphics[width=0.85\textwidth]{bab-05/gost-delapan-sbox}
  \caption{Tabel delapan Kotak-S ($S_1 \dots S_8$) $4 \times 4$ pada algoritma GOST.}
  \label{fig:gost-sbox}
\end{figure}

Keunikan terbesar GOST adalah bahwa \textbf{delapan kotak-S tidak ditetapkan secara kaku di dalam spesifikasi baku standar}. Kotak-S diperlakukan sebagai parameter kunci tambahan sekunder yang dapat disesuaikan untuk lembaga tertentu (misalnya bank sentral Rusia) atau bahkan dirahasiakan, memberikan tingkat kerahasiaan tambahan yang substansial.

\subsubsection{Penjadwalan Kunci (Key Schedule) GOST}

Kunci 256 bit dipecah secara langsung menjadi delapan buah subkunci 32-bit: $K_0, K_1, \dots, K_7$.
Urutan penggunaan subkunci selama 32 putaran adalah:
\begin{itemize}
  \item Putaran 1 sampai 24 (tiga siklus pertama): menggunakan urutan maju $K_0, K_1, \dots, K_7$ yang diulang sebanyak 3 kali.
  \item Putaran 25 sampai 32 (siklus terakhir): menggunakan urutan mundur $K_7, K_6, \dots, K_0$.
\end{itemize}
Dengan 32 putaran penuh dan kunci 256 bit, GOST sangat resisten terhadap serangan brute-force dan kriptanalisis diferensial.

\subsection{Algoritma RC5}

\textbf{RC5} adalah cipher blok simetris yang dirancang oleh Ronald L. Rivest pada tahun 1994. RC5 menjadi pelopor algoritma kriptografi modern yang bersifat \textbf{parametrik penuh}: ukuran kata, jumlah putaran, dan panjang kunci dapat dikonfigurasi secara bebas sesuai kebutuhan aplikasi.

Notasi formal algoritma dinyatakan sebagai $\mathbf{RC5-w/r/b}$:
\begin{itemize}
  \item $w$: Lebar word dalam bit ($w \in \{16, 32, 64\}$). Ukuran blok data adalah $2w$ bit (misalnya $2w = 64$ bit jika $w = 32$).
  \item $r$: Jumlah putaran ($r \in \{0, 1, \dots, 255\}$).
  \item $b$: Panjang kunci rahasia dalam byte ($b \in \{0, 1, \dots, 255\}$).
\end{itemize}
Konfigurasi standar yang paling luas diadopsi adalah $\mathbf{RC5-32/12/16}$ (blok 64 bit, 12 putaran, kunci 128 bit).

\begin{figure}[htbp]
  \centering
  \includegraphics[width=0.75\textwidth]{bab-05/rc5-round-diagram}
  \caption{Diagram alir komputasi satu putaran pada algoritma RC5 yang mengandalkan rotasi bergantung-data.}
  \label{fig:rc5-round}
\end{figure}

\subsubsection{Prinsip Perputaran Bergantung-Data (Data-Dependent Rotation)}

Inovasi utama dalam RC5 adalah penggunaan operasi \textbf{rotasi siklis yang besar pergeserannya ditentukan oleh nilai data itu sendiri} ($A \lll B$). Operasi ini menyuntikkan non-linearitas aljabar yang sangat kuat tanpa perlu mengalokasikan memori untuk tabel S-box, sehingga RC5 sangat kebal terhadap serangan waktu memori (\textit{timing attacks}).

Operasi dasar RC5 murni melibatkan tiga instruksi prosesor:
\begin{enumerate}
  \item Penjumlahan modulo $2^w$ (notasi $+$).
  \item Bitwise XOR (notasi $\oplus$).
  \item Rotasi bit siklis ke kiri yang bergantung nilai data: $A \lll B$ (di mana $A$ digeser sebanyak $B \bmod w$ posisi).
\end{enumerate}

\subsubsection{Algoritma Enkripsi RC5}

Blok plainteks dibagi menjadi dua kata $w$-bit, $A$ dan $B$. Larik subkunci diperluas dinamakan $S[0 .. 2r+1]$:
\begin{align}
  A &= A + S[0] \\
  B &= B + S[1] \\
  \text{Untuk } i &= 1 \text{ sampai } r \text{ lakukan:} \nonumber \\
  A &= ((A \oplus B) \lll B) + S[2i] \\
  B &= ((B \oplus A) \lll A) + S[2i + 1]
\end{align}
Keluaran akhir adalah gabungan kata $(A, B)$.

Pembangkitan subkunci $S$ menggunakan dua konstanta matematika ajaib berpresisi ganda: $P_w = \text{Odd}((e - 2) \cdot 2^w)$ (berbasis bilangan Euler $e$) dan $Q_w = \text{Odd}((\phi - 1) \cdot 2^w)$ (berbasis rasio emas $\phi$).

\subsection{Algoritma RC6}

\textbf{RC6} dirancang oleh Ronald Rivest, Matt Robshaw, Ray Sidney, dan Yiqun Lisa Yin sebagai proposal resmi untuk sayembara AES NIST. RC6 merupakan evolusi langsung dari RC5 yang disesuaikan untuk memenuhi persyaratan AES (blok 128 bit, kunci 128/192/256 bit, 20 putaran).

\begin{figure}[htbp]
  \centering
  \includegraphics[width=0.85\textwidth]{bab-05/rc6-feistel-diagram}
  \caption{Arsitektur empat register kata pada algoritma RC6.}
  \label{fig:rc6-feistel}
\end{figure}

\subsubsection{Perkalian Kuadratik Modular}

Untuk memproses blok 128 bit, RC6 membagi data menjadi empat register kata 32-bit: $(A, B, C, D)$. Keterbatasan RC5 adalah bahwa rotasi hanya ditentukan oleh beberapa bit terbawah dari data. Pada RC6, ditambahkan fungsi non-linear berbasis perkalian kuadratik modular:
\begin{equation}
  f(x) = (x \cdot (2x + 1)) \bmod 2^{32}
\end{equation}
Transformasi $f(x)$ memastikan bahwa seluruh bit masukan berkontribusi terhadap bit-bit penentu rotasi, melipatgandakan kecepatan difusi pada setiap putaran.

Satu putaran RC6 dirumuskan secara elegan:
\begin{align}
  t &= (B \cdot (2B + 1)) \lll 5 \\
  u &= (D \cdot (2D + 1)) \lll 5 \\
  A &= ((A \oplus t) \lll u) + S[2i] \\
  C &= ((C \oplus u) \lll t) + S[2i + 1] \\
  (A, B, C, D) &\leftarrow (B, C, D, A)
\end{align}
RC6 berhasil lolos sebagai salah satu dari lima finalis terbaik dalam sayembara AES.

\subsection{Algoritma Blowfish}

\textbf{Blowfish} dirancang oleh pakar keamanan terkemuka Bruce Schneier pada tahun 1993 sebagai cipher alternatif yang cepat, gratis, dan tidak terikat paten untuk menggantikan DES.

\begin{figure}[htbp]
  \centering
  \includegraphics[width=0.75\textwidth]{bab-05/blowfish-feistel-diagram}
  \caption{Diagram alir jaringan Feistel 16 putaran pada algoritma Blowfish.}
  \label{fig:blowfish-diagram}
\end{figure}

Karakteristik Blowfish:
\begin{itemize}
  \item \textbf{Ukuran Blok}: 64 bit.
  \item \textbf{Panjang Kunci}: Fleksibel, mulai dari 32 bit hingga \textbf{448 bit} (panjang standar 128 bit).
  \item \textbf{Struktur}: Jaringan Feistel 16 putaran dengan \textbf{P-array} berukuran 18 entri 32-bit ($P_1 \dots P_{18}$) dan \textbf{empat buah S-box besar} masing-masing berisi 256 entri 32-bit ($S_1 \dots S_4$).
\end{itemize}

\subsubsection{Inisialisasi S-box Berbasis Digit Pecahan $\pi$}

Salah satu aspek paling terkenal dari Blowfish adalah ketiadaan konstanta sembarang (\textit{no nothing-up-my-sleeve numbers}). Larik $P$ dan keempat kotak-S awalnya diinisialisasi secara deterministik menggunakan deretan digit heksadesimal dari bagian pecahan bilangan pi ($\pi$):
\begin{align}
  P_1 &= 243\text{F}6\text{A}88_{16}, \quad P_2 = 85\text{A}308\text{D}3_{16}, \quad \dots \\
  S_1[0] &= \text{D}1310\text{BA}6_{16}, \quad \dots
\end{align}

Selanjutnya, kunci masukan di-XOR ke dalam larik $P$, dan algoritma Blowfish digunakan untuk mengenkripsi deretan status berulang-ulang sebanyak 521 kali untuk menghasilkan kotak-S akhir yang sepenuhnya bergantung pada kunci (\textit{key-dependent S-boxes}).

Fungsi putaran Feistel $F(X_L)$ pada Blowfish didefinisikan sebagai:
\begin{equation}
  F(X_L) = ((S_1[a] + S_2[b] \bmod 2^{32}) \oplus S_3[c]) + S_4[d] \bmod 2^{32}
\end{equation}
di mana $a, b, c, d$ adalah empat byte yang menyusun kata 32-bit $X_L$.

Blowfish memiliki kecepatan komputasi software yang sangat impresif pada prosesor 32-bit dan hingga kini tidak ada serangan kriptanalisis praktis yang berhasil memecahkannya. Implementasinya masih digunakan secara luas dalam fungsi derivasi kunci dan hashing kata sandi \textbf{bcrypt} pada sistem operasi Unix/Linux.

\subsection{Finalis Sayembara AES: Twofish, Serpent, dan MARS}

Selain Rijndael (pemenang) dan RC6, tiga kandidat lainnya berhasil menembus babak final 5 besar sayembara AES:

\begin{figure}[htbp]
  \centering
  \begin{minipage}{0.48\textwidth}
    \centering
    \includegraphics[width=\textwidth]{bab-05/twofish-arsitektur}
    \caption{Arsitektur cipher Twofish dengan matriks MDS.}
    \label{fig:twofish-arsitektur}
  \end{minipage}
  \hfill
  \begin{minipage}{0.48\textwidth}
    \centering
    \includegraphics[width=\textwidth]{bab-05/serpent-spn-diagram}
    \caption{Struktur SPN 32 putaran pada algoritma Serpent.}
    \label{fig:serpent-spn}
  \end{minipage}
\end{figure}

\subsubsection{1. Twofish}
Dirancang oleh Bruce Schneier, John Kelsey, Doug Whiting, David Wagner, Chris Hall, dan Niels Ferguson. Twofish adalah penerus langsung Blowfish yang mengusung blok 128 bit dan kunci 128/192/256 bit. Twofish mempertahankan struktur Feistel 16 putaran dengan kotak-S dinamis yang bergantung kunci, serta menambahkan matriks difusi \textbf{MDS (Maximum Distance Separable)} $4 \times 4$ atas medan Galois $\mathrm{GF}(2^8)$ dan pembangkitan kunci berbasis Reed-Solomon. Twofish sangat fleksibel dan banyak digunakan dalam perangkat lunak enkripsi disk seperti VeraCrypt.

\subsubsection{2. Serpent}
Dirancang oleh Ross Anderson (Universitas Cambridge), Eli Biham (Technion), dan Lars Knudsen (DTU). Serpent mengadopsi arsitektur SPN 32 putaran dengan 32 kotak-S berukuran $4 \times 4$. Keunggulan revolusioner Serpent adalah implementasi \textbf{bitslicing}: seluruh 32 putaran dieksekusi menggunakan operasi logika boolean paralel (AND, OR, XOR, NOT) pada register prosesor 32-bit. Dalam sayembara AES, Serpent diakui secara bulat oleh para kriptografer dunia sebagai algoritma dengan \textbf{margin keamanan (security margin) paling tinggi dan paling konservatif}.

\subsubsection{3. MARS}
Dirancang oleh tim kriptografer elit IBM Research (dipimpin oleh Don Coppersmith, salah satu perancang asli DES). MARS mengusung arsitektur \textbf{heterogen multi-lapis}:
\begin{enumerate}
  \item \textbf{Forward Mixing}: 8 putaran pengacakan tanpa kunci untuk mencegah serangan penetrasi teks terang.
  \item \textbf{Cryptographic Core}: 16 putaran Feistel termodifikasi yang menggunakan perkalian data-dependent, penambahan subkunci, dan S-box $9 \times 32$ bit.
  \item \textbf{Backward Mixing}: 8 putaran pengacakan akhir tanpa kunci untuk mencegah serangan penetrasi teks sandi.
\end{enumerate}

\begin{figure}[htbp]
  \centering
  \includegraphics[width=0.85\textwidth]{bab-05/mars-mixing-diagram}
  \caption{Struktur arsitektur heterogen algoritma MARS: Forward Mixing, Cryptographic Core, dan Backward Mixing.}
  \label{fig:mars-diagram}
\end{figure}

\subsection{Penerapan Kriptografi Simetris dalam Kehidupan Sehari-hari}

Algoritma kriptografi simetris modern bekerja secara senyap di balik layar untuk melindungi hampir seluruh aktivitas digital masyarakat modern:

\subsubsection{1. Transaksi Finansial dan Mesin ATM (Automated Teller Machine)}

Setiap kali nasabah memasukkan kartu ATM dan mengetikkan kode \textit{Personal Identification Number} (PIN), keamanan dana nasabah dijamin oleh kriptografi simetris:
\begin{itemize}
  \item \textbf{Enkripsi Blok PIN}: Tombol papan ketik ATM (\textit{Encrypting PIN Pad} / EPP) langsung mengenkripsi PIN nasabah di dalam perangkat keras anti-sadap menjadi sebuah format terstandarisasi yang disebut \textbf{PIN Block} (ISO 9564 Format 0).
  \item \textbf{Peran HSM (Hardware Security Module)}: Enkripsi dan verifikasi PIN tidak pernah diproses oleh perangkat lunak biasa di server, melainkan dieksekusi di dalam modul kriptografi perangkat keras bersertifikasi FIPS 140-2 Level 3/4 yang disebut HSM.
  \item \textbf{Migrasi Algoritma}: Industri perbankan global berevolusi dari Single DES pada dekade 1980-an, beralih ke \textbf{Triple DES (TDES EDE2/EDE3)}, dan kini secara bertahap bermigrasi ke \textbf{AES-128 / AES-256}.
  \item \textbf{Otentikasi Pesan Transaksi}: Keaslian instruksi transfer dana dijamin menggunakan Message Authentication Code (MAC) simetris standar ANSI X9.9 dan ISO 9797-1.
\end{itemize}

\subsubsection{2. Komunikasi Bergerak dan Jaringan Telepon Seluler}

Dalam komunikasi telekomunikasi nirkabel, gelombang radio yang merambat bebas di udara sangat rentan disadap secara pasif oleh pihak mana pun yang memiliki antena penerima. Standar seluler mengandalkan kriptografi simetris pada setiap generasinya:
\begin{itemize}
  \item \textbf{Generasi 2G (GSM)}: Menggunakan cipher alir \textbf{A5/1} untuk mengenkripsi aliran percakapan suara nirkabel antara telepon genggam (\textit{Mobile Station}) dan stasiun pangkalan (\textit{Base Transceiver Station} / BTS).
  \item \textbf{Generasi 3G (UMTS)}: Menggantikan A5 dengan algoritma cipher blok \textbf{KASUMI} (varian dari MISTY1) dalam mode f8 (enkripsi) dan f9 (integritas data).
  \item \textbf{Generasi 4G (LTE)}: Menerapkan cipher alir modern \textbf{SNOW 3G} (standar EEA1/EIA1) dan \textbf{ZUC} (standar EEA3/EIA3), serta cipher blok \textbf{AES-128} dalam mode CTR dan CMAC (standar EEA2/EIA2).
  \item \textbf{Generasi 5G (5G NR)}: Memanfaatkan algoritma \textbf{AES-256} dan \textbf{ZUC-256} untuk menjamin kerahasiaan transmisi data ultra-cepat pada pita frekuensi milimeter.
\end{itemize}

\subsection{Tabel Komparasi Menyeluruh Block Cipher Simetris Modern}

\begin{figure}[htbp]
  \centering
  \includegraphics[width=0.9\textwidth]{bab-05/perbandingan-block-cipher-modern}
  \caption{Tabel komparasi karakteristik algoritma kriptografi simetri modern dari materi kuliah.}
  \label{fig:perbandingan-block-cipher}
\end{figure}

\begin{table}[htbp]
  \centering
  \small
  \resizebox{\textwidth}{!}{%
  \begin{tabular}{lcccccc}
    \toprule
    \textbf{Algoritma} & \textbf{Tahun} & \textbf{Perancang} & \textbf{Blok} & \textbf{Kunci (bit)} & \textbf{Rounds} & \textbf{Arsitektur} \\
    \midrule
    \textbf{DES} & 1977 & IBM / NBS & 64 & 56 & 16 & Feistel \\
    \textbf{3DES} & 1998 & ANSI / NIST & 64 & 112 / 168 & 48 & Feistel \\
    \textbf{GOST} & 1989 & Soviet Union & 64 & 256 & 32 & Feistel \\
    \textbf{Blowfish} & 1993 & Bruce Schneier & 64 & 32--448 & 16 & Feistel \\
    \textbf{RC5} & 1994 & Ronald Rivest & $2w$ & 0--2040 & 1--255 & Data-Dependent \\
    \textbf{RC6} & 1998 & Rivest dkk. & 128 & 128/192/256 & 20 & 4-Reg Feistel \\
    \textbf{MARS} & 1998 & IBM Research & 128 & 128--448 & 32 & Heterogen \\
    \textbf{Serpent} & 1998 & Anderson dkk. & 128 & 128/192/256 & 32 & SPN Bitslice \\
    \textbf{Twofish} & 1998 & Schneier dkk. & 128 & 128/192/256 & 16 & Feistel MDS \\
    \textbf{AES (Rijndael)} & 2001 & Daemen \& Rijmen & 128 & 128/192/256 & 10/12/14 & SPN $\mathrm{GF}(2^8)$ \\
    \bottomrule
  \end{tabular}%
  }
  \caption{Ringkasan komparasi arsitektur dan parameter algoritma block cipher terkemuka dunia.}
  \label{tab:komparasi-block-cipher-lengkap}
\end{table}

